\documentclass[11pt]{article}
\usepackage[margin=1in]{geometry}
\usepackage{amsmath,amssymb,amsthm,mathtools}
\usepackage[hidelinks]{hyperref}

\newtheorem{theorem}{Theorem}[section]
\newtheorem{lemma}[theorem]{Lemma}
\newtheorem{proposition}[theorem]{Proposition}
\newcommand{\tr}{\operatorname{tr}}
\newcommand{\Arg}{\operatorname{Arg}}
\newcommand{\one}{\mathbf 1}

\title{Positive Square Energy of Every Tricyclic Cactus}
\author{}
\date{25 July 2026}

\begin{document}
\maketitle

\begin{abstract}
We prove that every connected tricyclic cactus $G$ of order $n$ satisfies
$s^+(G)\geq n$.  The exact block-additive doubly nonnegative bound leaves
only the cycle triples $\{3,3,q\}$ with $q$ odd and $\{3,5,5\}$.  The first
family is settled by packing-two phase, induced cyclic territories, and a
product-subpartition comparison.  For $\{3,5,5\}$, a bridge separates a
one-cycle packet unless all cycles form one shared-cut cluster; the latter has
four cores, handled by an exact positive-coefficient phase certificate.
Arbitrary bridge blocks and arbitrary attached trees are allowed.
\end{abstract}

\section{Spectral and matching preliminaries}

All graphs are finite, simple, and undirected.  For the adjacency eigenvalues
$\lambda_1,\ldots,\lambda_n$, put
\[
 s^+(G)=\sum_{\lambda_j>0}\lambda_j^2,\quad
 s^-(G)=\sum_{\lambda_j<0}\lambda_j^2,\quad D(G)=s^+(G)-s^-(G).
\]
Since a connected tricyclic graph has $|E(G)|=n+2$,
\begin{equation}\label{eq:average}
 s^+(G)+s^-(G)=2n+4,\qquad s^+(G)=n+2+\frac12D(G).
\end{equation}
These square energies originate in spectral coloring bounds
\cite{AbiadEtAl}.  Liu--Tang--Zhang proved the general $n-1$ conjecture
\cite{LTZ}; the stronger edge-surplus conjecture of
Akbari--Kumar--Mohar--Pragada--Zhang asks for $s^+(G)\geq n$ when
$|E(G)|\geq n+1$ \cite{AKMPZ}.  Ning--Zeng proved the odd-unicyclic sign
theorem used below in its one-cycle form \cite{NingZeng}.

For positive activities $a=(a_v)$ define
\[
 Z_F(a)=\sum_{M\text{ matching of }F}\prod_{v\text{ unmatched by }M}a_v,
 \qquad Z_F(t)=Z_F((t)_v),
\]
and normalize the characteristic polynomial by
\[
 \Psi_G(t)=i^{-n}\det(itI-A(G))=\prod_j(t+i\lambda_j).
\]
Sachs' formula, grouped by collections $\mathcal C$ of pairwise
vertex-disjoint cycles, is
\begin{equation}\label{eq:sachs}
 \Psi_G(t)=\sum_{\mathcal C}\prod_{C\in\mathcal C}(-2i^{-|C|})
 Z_{G-V(\mathcal C)}(t).
\end{equation}
Thus a $3\pmod4$ cycle has multiplier $-2i$, and a $1\pmod4$ cycle has
multiplier $2i$.  If $\Theta_G(t)$ is the continuous argument of $\Psi_G(t)$
tending to zero at infinity, then
\begin{equation}\label{eq:coulson}
 D(G)=-\frac4\pi\int_0^\infty t\Theta_G(t)\,dt.
\end{equation}
Indeed, integrate
$\lambda|\lambda|=(2/\pi)\int_0^\infty
\lambda^3/(\lambda^2+t^2)\,dt$, sum over the spectrum, use
$\sum\lambda_j=0$, and integrate by parts.  This also fixes the sign and
branch in \eqref{eq:coulson}.

We use two standard exact reductions.  First, for every vertex partition
$(V_j)$, induced-subgraph superadditivity gives
\begin{equation}\label{eq:super}
 s^+(G)\geq\sum_j s^+(G[V_j])
\end{equation}
\cite{AKMP}.  Second, a tree branch oriented toward its attachment satisfies
\begin{equation}\label{eq:BP}
 Q_{u\to p}=\frac{Z_{T_{u\to p}}(t)}{Z_{T_{u\to p}-u}(t)}
 =t+\sum_{w\text{ child of }u}Q_{w\to u}^{-1}\geq t.
\end{equation}
Eliminating all branches therefore gives a common positive factor $K(t)$ and
activity $a_v=t+y_v$, $y_v\geq0$, at every retained vertex, in every Sachs
term.  This follows by splitting a matching according as its root is unmatched
or matched to one child.

\begin{lemma}[Phase and packet bounds]\label{lem:packets}
The following hold, with arbitrary tree attachments.
\begin{enumerate}
\item If the graph contains a cycle, every cycle is $3\pmod4$, and at most two are pairwise
vertex-disjoint, then $D>0$.
\item A unicyclic $C_q$ cactus with $q=1\pmod4$ satisfies
$s^+\geq h-\delta_q$, where $\delta_q=\sec(\pi/q)-1$.
\item A bicyclic cactus with blocks $C_3,C_q$, $q=1\pmod4$, satisfies
$s^+>h+1-\delta_q$.  Every bicyclic cactus satisfies $s^+\geq h$.
\end{enumerate}
\end{lemma}

\begin{proof}
In (1), \eqref{eq:sachs} has strictly negative imaginary part: singleton
cycle terms are negative imaginary, while empty and two-cycle terms are real.
Thus $\Theta<0$, and \eqref{eq:coulson} gives $D>0$.  This is the
packing-two theorem, including the unicyclic result of Ning--Zeng.

For (2), \eqref{eq:BP} gives
$\Psi/K=Z_{C_q}(a)+2i$.  Since $a_v\geq t$, coefficientwise monotonicity of
$Z$ bounds its phase by that of the bare cycle.  The eigenvalues
$2\cos(2\pi j/q)$ give
\begin{equation}\label{eq:cycleformula}
 s^+(C_q)=q+1-\sec(\pi/q),\qquad D(C_q)=-2\delta_q.
\end{equation}
Coulson integration proves (2).

For the mixed packet, retain the two cycles and their joining path.  Put
$A=Z_{H-V(C_3)}(a)$, $B=Z_{H-V(C_q)}(a)$, and let $D_0$ be the partition
after deleting both cycles when they are disjoint.  Sachs gives
\[
 \Psi/K=R+2i(B-A),\quad
 R=Z_H(a)\quad\hbox{or}\quad Z_H(a)+4D_0.
\]
Partitioning matchings by whether they use a boundary edge of $C_q$ gives
$Z_H=Z_{C_q}(a|_{C_q})B+E$, $E\geq0$.  Write
$Z_{C_q}(a|_{C_q})=Z_{C_q}(t)+L$, $L\geq0$.  Then
\[
 R-Z_{C_q}(t)(B-A)=E+LB+Z_{C_q}(t)A\ (+4D_0)>0.
\]
Hence $\Theta<\Arg(Z_{C_q}(t)+2i)$, so
$D>D(C_q)$ and \eqref{eq:average} in the bicyclic form proves (3)'s mixed
bound.  The final assertion is the established bicyclic-cactus theorem
\cite{AllBicyclic}; it follows there from the same sharp DNN reduction and
the two residual phase arguments.  These proofs use neither triangle-ear nor
edge-monotonicity.
\end{proof}

\section{The exact DNN reduction}

For a connected edge-containing graph define
\[
 \kappa(G)=\sup_{\substack{M\succeq0,\ M\geq0\\\one^TM\one>0}}
 \frac{4(\sum_{uv\in E(G)}\sqrt{M_{uv}})^2}{\one^TM\one}.
\]

\begin{lemma}[Sharp cactus DNN constant]\label{lem:dnn}
If a cactus has $b$ bridge blocks and cyclic blocks of lengths
$\ell_1,\ldots,\ell_r$, then
\[
 \kappa(G)=b+\sum_j\kappa(C_{\ell_j}),\qquad
 \kappa(C_\ell)=\begin{cases}\ell,&\ell\text{ even},\\
 \ell\sec^2(\pi/(2\ell)),&\ell\text{ odd},\end{cases}
\]
and $s^-(G)\leq\kappa(G)$.
\end{lemma}

\begin{proof}
Fold every nonedge entry $M_{uv}$ into the two diagonals by adding
$M_{uv}(e_u-e_v)(e_u-e_v)^T$.  The denominator and edge entries are
unchanged.  On $C_\ell$, rotational averaging and concavity of the square
root reduce to common edge entry $x$ and row sum $r$.  A Fourier eigenvalue
gives $r\geq4x$ for even $\ell$ and
$r\geq2(1+\cos(\pi/\ell))x$ for odd $\ell$.  Equality is attained by
$x(2cI+A(C_\ell))$, with $c=1$ or $c=\cos(\pi/\ell)$, proving the formula.
Also $\kappa(K_2)=1$.

For explicit cut-vector bookkeeping, suppose $G=G_1\vee_zG_2$.  In a folded
Gram system write $g_z=u+w+h$, where $u$ lies in the span $U$ of side 1,
$w$ lies in the orthogonal span $W$ of side 2, and $h\perp U+W$.  Use $u+h$
as the cut vector on side 1 and $w$ on side 2.  Edge inner products remain
unchanged and, if $a,b$ are the sums of the noncut vectors on the two sides,
\[
 \|a+b+u+w+h\|^2=\|a+u+h\|^2+\|b+w\|^2.
\]
Thus both the edge numerator sum and the denominator split.  Cauchy--Schwarz
gives $\kappa(G)\leq\kappa(G_1)+\kappa(G_2)$.  Conversely, glue near-optimal
Gram systems orthogonally and scale them so
$T_1/\kappa(G_1)=T_2/\kappa(G_2)$; this gives equality.  Iteration proves
block additivity.

Finally write $A=P-B$ for the positive and negative spectral parts.  The
Schur product theorem makes $M=B\circ B$ DNN, while
$s^-/2=-\sum_{uv\in E}B_{uv}\leq\sum_{uv\in E}|B_{uv}|$ and
$\one^TM\one=\tr B^2=s^-$.  The definition of $\kappa$ gives
$(s^-)^2\leq\kappa(G)s^-$, proving the spectral bound.
\end{proof}

Put $\epsilon_\ell=0$ for even $\ell$ and
$\epsilon_\ell=\ell\tan^2(\pi/(2\ell))$ for odd $\ell$.  Then
\begin{equation}\label{eq:eps}
 \epsilon_3=1,\quad \epsilon_5=5-2\sqrt5,\quad
 \epsilon_5+\epsilon_7<1,
\end{equation}
and $\epsilon_\ell$ strictly decreases through odd integers: after setting
$u=\pi/(2x)$ this is the monotonicity of $\tan^2u/u$, whose derivative is
positive because $2u>\sin u\cos u$.  The last inequality follows, for example,
from $\cos(\pi/7)>7/8$.

If the tricyclic cycle lengths are $p,q,r$, then
$b+p+q+r=n+2$.  Lemma~\ref{lem:dnn} yields
\begin{equation}\label{eq:dnn-bound}
 s^+(G)\geq n+2-\epsilon_p-\epsilon_q-\epsilon_r.
\end{equation}
Thus DNN proves $s^+\geq n$ unless the excess sum exceeds $2$.
Using \eqref{eq:eps} and monotonicity, this happens for exactly
\begin{equation}\label{eq:residual}
 \{3,3,q\}\ (q\geq3\text{ odd}),\qquad \{3,5,5\}.
\end{equation}
Indeed an even length contributes zero; without two triangles the only
possible maximum is $\epsilon_3+2\epsilon_5>2$, while replacing either five
by at least seven makes the sum $<2$.

\section{Two triangles and one odd cycle}

\begin{proposition}\label{prop:33q}
A cactus whose cyclic blocks are $C_3,C_3,C_q$, $q$ odd, satisfies $s^+>n$.
\end{proposition}

\begin{proof}
First let $q\equiv3\pmod4$; this includes the all-triangular case $q=3$.
If the cycle-packing number is at most two, Lemma~\ref{lem:packets}(1) gives
$D>0$ and hence $s^+>n+2$.  If the packing number is three, the cycles are
vertex-disjoint.  Cut bridge edges in the minimal block tree joining them and
assign every hanging tree to its attachment.  This partitions $V(G)$ into
three connected induced unicyclic territories.  Each has $D>0$, hence
$s^+$ greater than its order; \eqref{eq:super} gives $s^+(G)>n$.

Now let $q\equiv1\pmod4$ and $\delta_q=\sec(\pi/q)-1<1$.  If all cycles are
disjoint, choose an adjacent pair in the reduced three-node cycle tree.  Two
triangles form a packet with $s^+>h+1$, leaving a $C_q$ packet with
$s^+\geq h-\delta_q$; or a mixed pair has
$s^+>h+1-\delta_q$ and leaves a triangular packet with $s^+>h$.
Thus \eqref{eq:super} gives $s^+>n+1-\delta_q>n$.  The same bridge cut works
when one triangle is disjoint from the other triangle and $C_q$.

Suppose next that the triangles are disjoint and $C_q$ meets both, at distinct
cut vertices $x_1,x_2$.  Give $x_1$, its triangle, and branches off the
long-cycle side to one induced territory.  The complement is connected because
$C_q-x_1$ is a path.  Both territories are triangular unicyclic cacti, so
again $s^+>n$ by \eqref{eq:super}.

It remains that the triangles $T_1,T_2$ share a cut vertex.  Retain the cyclic
core and joining block paths and apply \eqref{eq:BP}.  For a cycle collection
$J$, put $D_J=Z_{H-V(J)}(a)$ when its cycles are disjoint and $D_J=0$
otherwise.  With $B=D_Q$ and $A_j=D_{T_j}$, Sachs gives
\[
 \Psi_G/K=R+iI,\quad
 R=Z_H+4D_{T_1,Q}+4D_{T_2,Q},\quad I=2(B-A_1-A_2).
\]
There is neither a double-triangle nor a triple term because the triangles
meet.  Partition matchings by boundary edges of $Q=C_q$:
\[
 Z_H=Z_Q(a|_Q)B+E,\quad E\geq0,\qquad
 Z_Q(a|_Q)=Z_q(t)+L,\quad L\geq0.
\]
Consequently
\[
 R-\frac{Z_q(t)I}{2}
 =E+LB+4D_{T_1,Q}+4D_{T_2,Q}+Z_q(t)(A_1+A_2)>0.
\]
Equivalently,
$\operatorname{Im}((\Psi_G/K)\overline{(Z_q+2i)})<0$.
The continuous arguments tend to zero at infinity, so they cannot cross and
$\Theta_G<\Arg(Z_q+2i)$.  Equations \eqref{eq:coulson}, \eqref{eq:cycleformula},
and \eqref{eq:average} give
$s^+(G)>n+2-\delta_q>n$.

For exhaustion: either the triangles meet, or they do not.  In the latter
case $Q$ meets both, at most one, or neither; these are respectively the
two-territory case, the separating-bridge case, and the reduced-tree case.
Thus every block incidence has been covered.
\end{proof}

\section{A triangle and two pentagons}

Join two cycle blocks in the \emph{shared-cut graph} when they share a cut
vertex.  If this graph is disconnected, contract each shared-cut component
and all acyclic paths between components.  The resulting cluster incidence is
a tree.  With three cycles, it has a leaf cluster consisting of exactly one
cycle: for a $2+1$ split the singleton is a leaf, and for a $1+1+1$ split any
leaf is.  Cutting the first bridge toward the rest partitions $G$ into two
connected induced territories, one containing exactly that cycle.

\begin{proposition}[Bridge-separable case]\label{prop:bridge}
If the shared-cut graph for cycle lengths $\{3,5,5\}$ is disconnected, then
$s^+(G)>n$.
\end{proposition}

\begin{proof}
Use the leaf-cluster cut above.  If the isolated packet is triangular, its
$s^+$ exceeds its order, while the remaining $\{5,5\}$ bicyclic packet has
$s^+$ at least its order by Lemma~\ref{lem:packets}(3).

If the isolated packet is pentagonal, put
$\delta_5=\sec(\pi/5)-1=\sqrt5-2$.  Lemma~\ref{lem:packets} gives
$s^+(U_5)\geq u-\delta_5$ and
$s^+(B_{3,5})>b+1-\delta_5$.  Superadditivity therefore gives
\[
 s^+(G)>n+1-2\delta_5=n+5-2\sqrt5>n.
\]
This includes three separate cyclic clusters; the cut is an actual bridge and
the two territories are induced, so no edge-monotonicity is invoked.
\end{proof}

\begin{proposition}[One shared-cut cluster]\label{prop:shared}
If the three blocks $C_3,C_5,C_5$ form one connected shared-cut cluster, then
\[
 s^+(G)>n+6-2\sqrt5>n.
\]
\end{proposition}

\begin{proof}
There are exactly four incidences up to isomorphism: a three-cycle bouquet;
the triangle as middle block with distinct cut vertices; or a pentagon as
middle block with its two cut vertices at cyclic distance one or two.  These
exhaust a connected three-node block tree, including coincident incidences.
All bridges are now branches off the cyclic core $H$, so \eqref{eq:BP} gives
independent positive core activities $a_v\geq t$ and a common positive factor.

Let $T,P,Q$ denote the triangle and pentagons and put
$D_J=Z_{H-V(J)}(a)$, interpreted as zero when the cycles in $J$ meet.  Sachs
gives $\Psi_G/K=R+iI$, where
\begin{align*}
 R&=Z_H+4D_{T,P}+4D_{T,Q}-4D_{P,Q},\\
 I&=-2D_T+2D_P+2D_Q.
\end{align*}
There is no triple term.  Let
$p=Z_P(a|_P)$ and $q=Z_Q(a|_Q)$ be the actual weighted isolated-pentagon
partitions.  Exact expansion over $\mathbb Z[a_v]$ gives the certificate
\begin{equation}\label{eq:phi}
 \Phi=2R(p+q)-I(pq-4)>0
\end{equation}
for every positive activity vector.  Table~\ref{tab:cert} records the complete
coefficient audit; the script constructs all matching partitions recursively,
checks the cores, and checks every coefficient.

\begin{table}[h]
\centering\small
\begin{tabular}{lrrr}
core & terms & min & max\\\hline
bouquet &2547&2&32\\
\multicolumn{4}{l}{\texttt{8112f2944f9823177afd48deccfb958ac960548e09d9d838e4965c33eb39e979}}\\[1mm]
triangle middle &2192&2&32\\
\multicolumn{4}{l}{\texttt{4fdb04cee38f0c0e2ac2de6dff7e641c2e190a3c018af8308a5c69e378de2ba2}}\\[1mm]
pentagon middle, distance 1&2925&2&36\\
\multicolumn{4}{l}{\texttt{bfa73346f169f28ec6109418ce22fe44daaae6b081ade30074932b139f0828f4}}\\[1mm]
pentagon middle, distance 2&2895&2&36\\
\multicolumn{4}{l}{\texttt{5c6c213471a856cb743afda8e62407d7d27de5984b612c70ab411499011db437}}\\
\end{tabular}
\caption{Exact coefficient certificates for $\Phi$.}\label{tab:cert}
\end{table}

Now $(p+2i)(q+2i)=(pq-4)+2i(p+q)$, and \eqref{eq:phi} is precisely
\[
 \operatorname{Im}((p+2i)(q+2i)\overline{(R+iI)})>0.
\]
The continuous arguments all tend to zero at infinity.  Strict positivity
prevents a branch crossing, hence
\[
 \Theta_G(t)<\Arg(p+2i)+\Arg(q+2i).
\]
Because $a_v\geq t$, each weighted isolated-pentagon phase is at most the bare
$C_5$ phase.  Coulson integration and
$D(C_5)=4-2\sqrt5$ therefore give
$D(G)>2D(C_5)=8-4\sqrt5$.  Equation \eqref{eq:average} proves the claim.
\end{proof}

\section{Conclusion}

\begin{theorem}\label{thm:main}
Every connected tricyclic cactus $G$ on $n$ vertices satisfies
\[
 \boxed{s^+(G)\geq n}.
\]
\end{theorem}

\begin{proof}
The sharp DNN bound \eqref{eq:dnn-bound} proves the assertion unless the cycle
length triple is one of \eqref{eq:residual}.  Proposition~\ref{prop:33q}
settles every $\{3,3,q\}$.  For $\{3,5,5\}$ the shared-cut graph is either
disconnected, covered by Proposition~\ref{prop:bridge}, or connected, covered
by Proposition~\ref{prop:shared}.  These alternatives exhaust all cycle
lengths, all block-cut incidences, all bridges, and all attached trees.
\end{proof}

\appendix
\section{Exact reproduction and local proof objects}

From the repository root run
\begin{verbatim}
python positive-square-energy/experiments/c3_c5_c5_shared_cluster_certificate.py
\end{verbatim}
to reproduce Table~\ref{tab:cert}.  The public local manuscripts used for the
packet lemmas and audited assembly are
\begin{verbatim}
all-bicyclic-cacti/paper.tex
packing-two-square-energy/paper.tex
research/c3-c3-cq-tricyclic-cactus-2026-07-25.md
research/c3-c5-c5-shared-cluster-2026-07-25.md
\end{verbatim}

\begin{thebibliography}{99}
\bibitem{AbiadEtAl}
A.~Abiad, L.~de Lima, D.~N.~Desai, K.~Guo, L.~Hogben, and J.~Madrid,
\emph{Positive and negative square energies of graphs},
Electron. J. Linear Algebra 39 (2023), 307--326.

\bibitem{AKMP}
S.~Akbari, H.~Kumar, B.~Mohar, and S.~Pragada,
\emph{A linear lower bound for the square energy of graphs},
Electron. J. Combin. 32(3) (2025), Paper No. P3.53.

\bibitem{AKMPZ}
S.~Akbari, H.~Kumar, B.~Mohar, S.~Pragada, and S.~Zhang,
\emph{Refinement of a conjecture on positive square energy of graphs},
arXiv:2506.07264, 2025.

\bibitem{LTZ}
Y.~Liu, Q.~Tang, and S.~Zhang,
\emph{The positive and negative square-energy conjecture},
arXiv:2607.18031, 2026.

\bibitem{NingZeng}
B.~Ning and J.~Zeng,
\emph{A proof of a conjecture on positive and negative square energies of
unicyclic graphs}, arXiv:2605.24668, 2026.

\bibitem{AllBicyclic}
Companion manuscript, \emph{Positive Square Energy of Every Bicyclic Cactus},
2026; repository path \texttt{all-bicyclic-cacti/paper.tex}.
\end{thebibliography}
\end{document}
